\newpage
\section{Introdução}
\label{sc:introducao}

O avanço contemporâneo dos Modelos de Linguagem de Grande Porte (LLMs) expandiu significativamente as fronteiras da Inteligência Artificial, permitindo a execução de tarefas complexas que variam desde a análise de documentos legais até a geração automatizada de códigos de programação. Contudo, apesar do notável desempenho em tarefas estritamente textuais, os LLMs tradicionais enfrentam limitações severas ao lidar com conhecimentos em constante evolução e dados de domínios específicos, frequentemente resultando em respostas alucinatórias ou desatualizadas. Para mitigar esse problema, a arquitetura de Geração Aumentada de Recuperação (RAG) consolidou-se como uma abordagem não-paramétrica indispensável \cite{andriopoulos2023augmenting}, integrando a capacidade gerativa dos modelos à precisão de bases de conhecimento externas indexadas em bancos de dados vetoriais \cite{ma2024vectorsurvey}. No entanto, a implementação prática de novos sistemas RAG ainda introduz desafios complexos de engenharia e infraestrutura que exigem investigação, tais como restrições de tempo de execução, saturação de limites de tokens e o consumo elevado de recursos de hardware.  

Diante desse cenário, este trabalho propõe-se a realizar uma análise investigativa das metodologias de otimização aplicadas aos processos de RAG. O escopo delineado para a pesquisa abrangerá uma avaliação comparativa detalhada entre abordagens clássicas e avançadas de recuperação — especificamente os métodos Stuff, Refine, Map Reduce, Map Re-rank, Query Step-Down e Reciprocal RAG \cite{topsakal2023langchain}. Pretende-se analisar os impactos diretos que a escolha dessas técnicas exerce sobre três pilares fundamentais de desempenho no ecossistema de inteligência artificial: a acurácia semântica das respostas geradas (mensurada por métricas de similaridade de cosseno), a eficiência no consumo de tokens e a utilização de infraestrutura de hardware, envolvendo o consumo de CPU, memória RAM e tempo de resposta. Esse recorte metodológico justifica-se pela necessidade premente de compreender os múltiplos tradeoffs técnicos existentes entre a precisão da resposta entregue ao usuário e os custos operacionais de computação em nuvem em ambientes de alta escala.  

O objetivo principal desta proposta de pesquisa é projetar e sintetizar um panorama analítico sobre a eficiência de arquiteturas RAG, buscando identificar as melhores configurações de componentes para diferentes cenários de aplicação. Especificamente, o trabalho investigará como a escolha combinada de bancos de dados vetoriais (como ChromaDB, FAISS e Pinecone), modelos de embedding (como Bge-en-small, Text-embedding-v3-large e Cohere-en-v3) e filtros de compressão contextual (limiares de similaridade e redundância) afetam o equilíbrio entre o gerenciamento de recursos e a qualidade do texto final gerado pelo modelo. Para dar sustentação teórica e metodológica a este projeto, a pesquisa tomará como base empírica e ponto de partida os dados e parâmetros discutidos na literatura recente sobre otimização sistemática por varredura em grade (grid-search), o que inclui as metodologias de testes cruzados em múltiplos domínios textuais, tais como relatórios financeiros corporativos (SEC 10Q), documentação científica (Llama2 ArXiv) e exames da área médica (MedQA) \cite{sakar2025rag}.  

A revisão de literatura que fundamentará este projeto terá como foco publicações recentes indexadas em bases de prestígio acadêmico internacional, com ênfase em termos-chave estruturantes como Retrieval-Augmented Generation, Vector Databases, Contextual Compression e Embedding Models. Por fim, a proposta de desenvolvimento deste estudo está estruturada de forma a aprofundar essas discussões técnicas nas seções subsequentes. O plano de trabalho prevê inicialmente a apresentação do referencial teórico acerca do funcionamento matemático da similaridade de cosseno e dos algoritmos de busca por vizinhos mais próximos. Em seguida, detalha-se o cronograma e a metodologia planejada para a análise comparativa dos indicadores de desempenho de acurácia, consumo de hardware e eficácia dos filtros de compressão. Por fim, serão delineados os resultados esperados e as contribuições pretendidas para o desenvolvimento de aplicações práticas baseadas nas conclusões a serem estabelecidas pelo estudo. 